\chapter{Spazi vettoriali}

\thispagestyle{empty}

Il cuore del corso di Geometria e Algebra Lineare è sicuramente lo studio degli spazi vettoriali. In questa sezione vedremo le definizioni, le proprietà e i teoriemi fondamentali che riguardano gli spazi vettoriali, i sottospazi, il concetto di base e di dimensione. \\

\section{Definizioni e proprietà}
La definizione di spazio vettoriale è strettamente legata alla definizione di campo, che può essere così definito
\begin{definition}[Campo]
    \label{def:campo}
    Un \textbf{campo} è un insieme $\mathbb{K}$ dotato di due operazioni $+_\mathbb{K}: \mathbb{K} \times \mathbb{K} \to \mathbb{K}$ e $\cdot_\mathbb{K}: \mathbb{K} \times \mathbb{K} \to \mathbb{K}$ tali che:
    \begin{itemize}
        \item $(\mathbb{K}, +)$ è un gruppo abeliano, ovvero l'operazione $+$ gode delle seguenti proprietà:
        \begin{enumerate}[label=\protect\circled{\arabic*}]
            \item $\forall a, b, c \in \mathbb{K}, (a +_\mathbb{K} b) +_\mathbb{K} c = a +_\mathbb{K} (b +_\mathbb{K} c)$ (proprietà associativa);
            \item $\forall a, b \in \mathbb{K}, a +_\mathbb{K} b = b +_\mathbb{K} a$ (proprietà commutativa);
            \item $\exists 0_\mathbb{K} \in \mathbb{K} : \forall a \in \mathbb{K}, 0_\mathbb{K} +_\mathbb{K} a = a +_\mathbb{K} 0_\mathbb{K} = a$ (esistenza dell'elemento neutro);
            \item $\forall a \in \mathbb{K}, \exists (-a) \in \mathbb{K} : a +_\mathbb{K} (-a) = (-a) +_\mathbb{K} a = 0$ (esistenza dell'elemento inverso).
        \end{enumerate}
        \item $(\mathbb{K} \setminus \{ 0_\mathbb{K} \}, \cdot)$ è un gruppo abeliano, ovvero l'operazione $\cdot$ gode delle seguenti proprietà:
        \begin{enumerate}[label=\protect\circled{\arabic*}]
            \item $\forall a, b, c \in \mathbb{K}, (a \cdot_\mathbb{K} b) \cdot_\mathbb{K} c = a \cdot_\mathbb{K} (b \cdot_\mathbb{K} c)$ (proprietà associativa);
            \item $\forall a, b \in \mathbb{K}, a \cdot_\mathbb{K} b = b \cdot_\mathbb{K} a$ (proprietà commutativa);
            \item $\exists \mathds{1}_{\mathbb{K}} \in \mathbb{K} : \mathds{1}_{\mathbb{K}} \cdot_\mathbb{K} a = a \cdot_\mathbb{K} \mathds{1}_{\mathbb{K}} = a$ (esistenza dell'elemento neutro);
            \item $\forall a \in \mathbb{K} \setminus \{ 0_\mathbb{K} \}, \exists a^{-1} \in \mathbb{K} : a \cdot_\mathbb{K} a^{-1} = a^{-1} \cdot_\mathbb{K} a = 1$ (esistenza dell'elemento inverso).
        \end{enumerate}
    \end{itemize}
\end{definition}
\begin{remark}
    Fare attenzione ai pedici sulle operazioni e sugli elementi: infatti, nel corso di questi appunti avremo a che fare spesso con delle operazioni che sono definite su insiemi diversi ma che vengono preservate dai morfismi; tuttavia, per evitare confusione, cercherò il più possibile di denotare le operazioni e gli elementi con pedici che tolgano ogni dubbio sull'insieme di appartenenza.
\end{remark}
\begin{example}
    L'esempio più lampante e semplice di campo è chiaramente l'insieme dei numeri razionali $\mathbb{Q}$ e l'insieme dei numeri reali $\mathbb{R}$ (e chiaramente anche il campo dei numeri complessi $\mathbb{C}$). Un esempio di campo meno noto è sicuramente $\mathbb{Z}_2$, ovvero l'insieme dei numeri interi modulo $2$, che contiene solo due elementi, $0$ e $1$, e le operazioni di somma e prodotto sono definite come segue:
    \begin{itemize}
        \item $0 + 0 = 0$, $0 + 1 = 1$, $1 + 0 = 1$, $1 + 1 = 0$;
        \item $0 \cdot 0 = 0$, $0 \cdot 1 = 0$, $1 \cdot 0 = 0$, $1 \cdot 1 = 1$.
    \end{itemize}
    e si può verificare che tale insieme è effettivamente un campo, in quanto le due operazioni possiedono tutte le proprietà richieste dalla definizione data in \ref{def:campo}.
\end{example}
\begin{definition}[Spazio vettoriale]
    \label{def:sp_vett} Sia dato un insieme $V$, un campo $\mathbb{K}$ e siano date due operazioni $+: V \times V \to V$ e $\cdot: \mathbb{K} \times V \to V$. Diremo che $(V, \mathbb{K}, +, \cdot)$ è un $\mathbb{K}-$spazio vettoriale (o più semplicemente uno spazio vettoriale sul campo $\mathbb{K}$) se le due operazioni soddisfano le seguenti proprietà:
    \begin{itemize}
        \item $(V, +)$ è un gruppo abeliano, ovvero l'operazione $+$ gode delle seguenti proprietà:
        \begin{enumerate}[label=\protect\circled{\arabic*}]
            \item $\forall \underline{u}, \underline{v}, \underline{z} \in V, (\underline{u} + \underline{v}) + \underline{z} = \underline{u} + (\underline{v} + \underline{z})$ (proprietà associativa);
            \item $\forall \underline{v}, \underline{w} \in V, \underline{v} + \underline{w} = \underline{w} + \underline{v}$ (proprietà commutativa);
            \item $\exists \underline{0} \in V : \forall \underline{v} \in V, \underline{0} + \underline{v} = \underline{v} + \underline{0} = \underline{v}$ (esistenza dell'elemento neutro);
            \item $\forall \underline{v} \in V, \exists (-\underline{v}) \in V : \underline{v} + (-\underline{v}) = (-\underline{v}) + \underline{v} = \underline{0}$ (esistenza dell'elemento inverso).
        \end{enumerate}
        \item l'operazione $\cdot: \mathbb{K} \times V \to V$ gode delle seguenti proprietà:
        \begin{enumerate}[label=\protect\circled{\arabic*}]
            \item $\forall \alpha \in \mathbb{K}, \forall \underline{v}, \underline{w} \in V, \alpha \cdot (\underline{v} + \underline{w}) = \alpha \cdot \underline{v} + \alpha \cdot \underline{w}$ (proprietà distributiva rispetto alla $+$);
            \item $\forall \alpha, \beta \in \mathbb{K}, \forall \underline{v} \in V, (\alpha +_{\mathbb{K}} \beta) \cdot \underline{v} = \alpha \cdot \underline{v} + \beta \cdot \underline{v}$ (proprietà distributiva rispetto alla $+_\mathbb{K}$);
            \item $\forall \alpha, \beta \in \mathbb{K}, \forall \underline{v} \in V, (\alpha \beta) \cdot \underline{v} = \alpha \cdot (\beta \underline{v})$ (proprietà associativa della $\cdot$);
            \item $\exists \mathds{1} \in \mathbb{K} : \forall \underline{v} \in V, \mathds{1} \cdot \underline{v} = \underline{v}$ (esistenza dell'elemento neutro per $\cdot$).
        \end{enumerate}
        Le due operazioni $+$ e $\cdot$ vengono dette, rispettivamente, \textbf{somma} (o legge di composizione interna) e \textbf{prodotto per scalare} (o legge di composizione esterna).
    \end{itemize}
\end{definition}
\begin{remark}
    Concettualmente, sebbene siano simili, la definizione di spazio vettoriale è diversa da quella di campo: infatti, sul campo abbiamo due operazioni che sono definite sul prodotto cartesiano del campo con sé stesso e restituiscono un elemento del campo; mentre nel caso dello spazio vettoriale si ha un insieme $V$ su cui sono definite due operazioni, una che è definita sul prodotto cartesiano di $V$ con sé stesso e restituisce un elemento di $V$, e l'altra che è definita sul prodotto cartesiano del campo con l'insieme $V$ e restituisce un elemento di $V$.
\end{remark}
\begin{remark}
    Si noti che uno spazio vettoriale è definito completamente solamente se definiamo le due operazioni $+$ e $\cdot$, oltre al campo. E' possibile, come vedremo in un esempio, andare a fornire la struttura di spazio vettoriale ad un medesimo insieme definendo le due operazioni in modo opportuno.
\end{remark}
Tramite questa definizione, possiamo andare a mostrare alcuni risultati validi nei campi che sono altrettanto validi negli spazi vettoriali.
\begin{prop}[l'elemento neutro rispetto a $+$ è unico]
    Sia $(V, \mathbb{K}, +, \cdot)$ uno spazio vettoriale sul campo $\mathbb{K}$. Allora l'elemento neutro rispetto all'operazione $+$ è unico, ovvero
    $$
    \exists! \underline{0} \in V : \forall \underline{v} \in V, \underline{0} + \underline{v} = \underline{v} + \underline{0} = \underline{v}.
    $$
\end{prop}
\begin{prop}
    \label{prop:0v_equal_0} Sia $(V, \mathbb{K}, +, \cdot)$ uno spazio vettoriale sul campo $\mathbb{K}$. Allora
    $$
    \forall \underline{v} \in V, 0_{\mathbb{K}} \cdot \underline{v} = \underline{0}.
    $$
\end{prop}
\begin{prop}
    Sia $(V, \mathbb{K}, +, \cdot)$ uno spazio vettoriale sul campo $\mathbb{K}$. Allora
    $$
    \forall \underline{v} \in V, (-1_\mathbb{K}) \cdot \underline{v} = (-\underline{v}).
    $$
\end{prop}
\begin{exercise}
    \label{ex:1.1} Mostrare le tre proposizioni precedenti. (\emph{Suggerimento}: per la prima si può procedere per assurdo, negando la tesi; per la seconda si può mostrare direttamente usando il fatto che $0_\mathbb{K} = 0_\mathbb{K} +_\mathbb{K} 0_\mathbb{K}$ e la terza si può mostrare direttamente usando la seconda.)
\end{exercise}
Vediamo adesso qualche esempio di spazio vettoriale per fissare meglio i concetti.
\begin{example}[$\mathbb{R}^2$ è uno spazio vettoriale sul campo $\mathbb{R}$]
    Consideriamo l'insieme $\mathbb{R}^2 = \mathbb{R} \times \mathbb{R}$: gli elementi di questo insieme sono coppie ordinate di numeri reali, che possiamo denotare come $\underline{v} = (v_1, v_2)$ con $v_1, v_2 \in \mathbb{R}$. Possiamo definire le due operazioni $+$ e $\cdot$ come segue:
    \begin{align*}
        &\underline{v} + \underline{w} = \begin{pmatrix} v_1 \\ v_2 \end{pmatrix} + \begin{pmatrix} w_1 \\ w_2 \end{pmatrix} = \begin{pmatrix} v_1 + w_1 \\ v_2 + w_2 \end{pmatrix}
        &\alpha \cdot \underline{v} = \alpha \cdot \begin{pmatrix} v_1 \\ v_2 \end{pmatrix} = \begin{pmatrix} \alpha v_1 \\ \alpha v_2 \end{pmatrix}
    \end{align*}
    ed è possibile verificare che queste due operazioni soddisfano sempre le proprietà richieste dalla definizione \ref{def:sp_vett}
\end{example}
\begin{example}[$\mathbb{R}^n$ è uno spazio vettoriale sul campo $\mathbb{R}$]
    Partendo dall'idea dell'esempio precedente, possiamo iterare il ragionamento e definire l'insieme $\mathbb{R}^n = \underbrace{\mathbb{R} \times \ldots \times \mathbb{R}}_{\text{n volte}}$: gli elementi di questo insieme sono $n$-tuple ordinate di numeri reali, che possiamo denotare come $\underline{v} = (v_1, \ldots, v_n)$ con $v_i \in \mathbb{R} \forall i \in \{ 1, \ldots, n \}$. Possiamo definire la somma e il prodotto per scalare come fatto prima, sommando componente per componente e moltiplicando componente per componente:
    \begin{align*}
        &\underline{v} + \underline{w} = \begin{pmatrix} v_1 \\ \vdots \\ v_n \end{pmatrix} + \begin{pmatrix} w_1 \\ \vdots \\ w_n \end{pmatrix} = \begin{pmatrix} v_1 + w_1 \\ \vdots \\ v_n + w_n \end{pmatrix} \\
        &\alpha \cdot \underline{v} = \alpha \cdot \begin{pmatrix} v_1 \\ \vdots \\ v_n \end{pmatrix} = \begin{pmatrix} \alpha v_1 \\ \vdots \\ \alpha v_n \end{pmatrix}.
    \end{align*}
\end{example}
\begin{example}[$\mathbb{C}^n$ è uno spazio vettoriale sul campo $\mathbb{C}$]
    Analogamente a quanto fatto per $\mathbb{R}^n$, possiamo riadattare il ragionamento fatto per $\mathbb{R}^n$ ammettendo che le componenti delle $n$-tuple ordinate siano numeri complessi, ovvero $\mathbb{C}^n \ni \underline{v} = \begin{pmatrix} v_1 \\ \vdots \\ v_n \end{pmatrix}$ con $v_i \in \mathbb{C} \forall i \in \{ 1, \ldots, n \}$.
\end{example}
\begin{example}[le soluzioni di un'equazione differenziale lineare sono uno spazio vettoriale]
    Sia data un'equazione differenziale lineare di ordine $n$ del tipo
    $$
        \sum_i^n a_i y^{(i)}(x) = 0 \text{ con } a_i \in \mathbb{R}
    $$
    allora possiamo mostrare che l'insieme delle soluzioni di questa equazione, denotato con $V$ gode della struttura di spazio vettoriale sul campo $\mathbb{R}$ definendo le due operazioni come segue: prese due soluzioni $y_1, y_2 \in V$ allora poniamo
    \begin{align*}
        &(y_1 + y_2)(x) = y_1(x) + y_2(x) \, \, \forall x \in \mathbb{R}
        &(\alpha y_1)(x) = \alpha y_1(x) \, \, \forall x \in \mathbb{R}
    \end{align*}
    e osserviamo che queste due operazioni sono chiuse, ovvero restituiscono un elemento che è ancora una soluzione dell'equazione differenziale lineare. Infatti
    $$
        \sum_i^n a_i (y_1 + y_2)^{(i)}(x) = \sum_i^n a_i y_1^{(i)} + \sum_i^n a_i y_2^{(i)} = 0 + 0 = 0,
    $$
    dove abbiamo usato il fatto che l'operazione di derivata è distributiva rispetto alla somma e il fatto che $y_1, y_2$ fossero soluzioni dell'equazione differenziale lineare, ovvero $\sum_i^n a_i y_j^{(i)} = 0$, con $j=1, 2$. Similmente, per quanto riguarda il prodotto per scalare, si usa il fatto che le costanti non sono coinvolte nel processo di derivazione, dunque
    $$
        \sum_i^n a_i(\alpha y_1)^{(i)}(x) = \sum_i^n a_i \alpha y_1^{(i)}(x) = \alpha \sum_i^n a_i y_1^{(i)}(x) = 0,
    $$
    dove si è sempre usato il fatto che $y_1$ fosse una soluzione dell'equazione differenziale lineare. Le proprietà richieste dalla definizione \ref{def:sp_vett} sono banalmente verificabili.
\end{example}
\begin{example}[spazio dei polinomi a coefficienti reali di grado $\leq n$ è uno spazio vettoriale]
    Sia dato l'insieme dei polinomi a coefficienti reali con grado $\leq n$ $\mathbb{R}_{\leq n}[x]$ su $\mathbb{R}$, allora possiamo definire le due operazioni come segue: presi due polinomi $p_1, p_2 \in \mathbb{R}_{\leq n}[x]$ possiamo definire il seguente polinomio
    \begin{align*}
    &(p_1 + p_2)(x) = p_1(x) + p_2(x) \, \, \forall x \in \mathbb{R}
    &(\alpha p_1)(x) = \alpha p_1(x) \, \, \forall x \in \mathbb{R}.
    \end{align*}
    Si può verificare che queste due operazioni sono chiuse, ovvero restituiscono sempre un polinomio a coefficienti reali di grado $\leq n$: infatti se $p_1(x) = \sum\limits_{i=1}^n a_i x^i$ e $p_2(x) = \sum\limits_{i=1}^n b_i x^i$ con $a_i, b_i \in \mathbb{R}$ allora
    \begin{align*}
        &(p_1 + p_2)(x) = \sum_{i=1}^n a_i x^i + \sum_{i=1}^n b_i x^i = \sum_{i=1}^n (a_i + b_i) x^i \\
        &(\alpha p_1)(x) = \alpha \sum_{i=1}^n a_i x^i = \sum_{i=1}^n (\alpha a_i) x^i
    \end{align*}
\end{example}
\begin{example}[Uno spazio vettoriale è definito solamente conoscendo la legge interna ed esterna]
    Fissiamo come spazio vettoriale $V = \mathbb{R}^n$, che è chiaramente uno spazio vettoriale sul campo $\mathbb{R}$, e definiamo le due operazioni $\oplus$ e $\odot$ come
    \begin{align*}
        &\underline{v} \oplus \underline{w} = \begin{pmatrix} v_1 \\ \vdots \\ v_n \end{pmatrix} \oplus \begin{pmatrix} w_1 \\ \vdots \\ w_n \end{pmatrix} = \begin{pmatrix} v_1 + w_1 + 2 \\ \vdots \\ v_n + w_n +2 \end{pmatrix} \\
        &\alpha \odot \underline{v} = \alpha \odot \begin{pmatrix} v_1, \ldots, v_n \end{pmatrix} = \begin{pmatrix} \alpha v_1 - 2(1-\alpha) \\ \alpha v_2 - 2(1-\alpha) \\ \vdots \\ \alpha v_n - 2(1-\alpha) \end{pmatrix}.
    \end{align*}
    Si può verificare che queste due operazioni sono chiuse, ovvero restituiscono sempre un elemento di $V$, e soddisfano le proprietà richieste dalla definizione \ref{def:sp_vett} di spazio vettoriale. Ne deduciamo che uno spazio vettoriale è completamente determinato \textbf{solo} conoscendo le due operazioni $+: V \times V \to V$ e $\cdot: \mathbb{K} \times V \to V$, sebbene in questo spazio vettoriale non è il vettore $(0, \ldots, 0)$ ad essere l'elemento neutro rispetto all'operazione $+$, ma il vettore $(-2, \ldots, -2)$; tuttavia, è interessante notare come
    comunque valgano le proprietà dimostrate nella proposizione \ref{prop:0v_equal_0} (infatti, $0 \odot \underline{v} = \begin{pmatrix} 0 v_1 - 2(1-0) \\ \vdots \\ 0v_n - 2(1-0) \end{pmatrix} = \begin{pmatrix} -2 \\ \vdots \\ -2 \end{pmatrix} )$ e nelle due successivi mostrate.
\end{example}
\section{Sottospazi}

Siamo ora interessati a studiare dei sottoinsiemi di uno spazio vettoriale che siano a loro volta degli spazi vettoriali rispetto all'operazione di somma e di prodotto per scalare. I sottospazi più comuni che possiamo immaginare, come vedremo, sono le rette e i piani che abbiamo in $\mathbb{R}^3$ (le rette anche in $\mathbb{R}^2$): la potenza dell'approccio che svilupperemo è il fatto che potremo trattare oggetti, come gli iperpiani, come se fossero oggetti i nostri usuali piani. 
\begin{definition}[sottospazio]
    Sia $(V, \mathbb{K}, +, \cdot)$ uno spazio vettoriale sul campo $\mathbb{K}$ e sia $W \subseteq V$ un sottoinsieme di $V$. Diremo che $W$ è un \textbf{sottospazio vettoriale} se esso è uno spazio vettoriale sul campo $\mathbb{K}$ rispetto alle operazioni $+$ e $\cdot$ definite su $V$.
\end{definition}
\begin{remark}
    Ogni spazio vettoriale è un sottospazio di sé stesso, siccome è chiaro che $V \subseteq V$.
\end{remark}
\begin{prop}
    \label{prop:ssp_vett_iff}
    Sia $(V, \mathbb{K}, +, \cdot)$ uno spazio vettoriale sul campo $\mathbb{K}$. Allora $W \subseteq V$ è un sottospazio vettoriale se e solo se
    \begin{enumerate}[label=\protect\circled{\arabic*}]
        \item $\forall \underline{v}, \underline{w} \in V, (\underline{v}, \underline{w} \in W) \implies \underline{v} + \underline{w} \in W$;
        \item $\forall \alpha \in \mathbb{K}, \forall \underline{v} \in V, (\underline{v} \in W) \implies (\alpha \underline{v} \in W)$.
    \end{enumerate}
\end{prop}
\begin{proof} Mostriamo le due implicazioni. \\
    $\boxed{\Rightarrow}$: se $W \subseteq V$ è un sottospazio vettoriale, allora per definizione esso è uno spazio vettoriale sul campo $\mathbb{K}$, dunque le due implicazioni
    \begin{align*}
        &\underline{v}, \underline{w} \in W \implies \underline{v} + \underline{w} \in W;
        &\underline{v} \in W \implies \alpha \underline{v} \in W;
    \end{align*}
    sono sempre vere. \\
    $\boxed{\Leftarrow}$: se valgono le due implicazioni, allora è chiaro che le due operazioni $+$ e $\cdot$ sono chiuse su $W$, dunque è sufficiente mostrare che valgono le proprietà richiesta dalla definizione \ref{def:sp_vett} di spazio vettoriale. Tuttavia, siccome queste valgono per ogni elemento di $V$, è chiaro che varranno anche per ogni elemento di $W$ siccome $W \subseteq V$; dunque possiamo concludere che $W$ è un sottospazio vettoriale di $V$.  
\end{proof}
\textbf{Attenzione}: da adesso in poi, per comodità di scrittura, non scriverò più esplicitamente l'operazione $\cdot$ per indicare il prodotto per scalare, ma mi limiterò a scrivere direttamente $\alpha \underline{v}$ sottointendendo il prodotto per scalare.
\begin{definition}
    Sia $(V, \mathbb{K}, +, \cdot)$ uno spazio vettoriale sul campo $\mathbb{K}$ e siano $\underline{v_1}, \ldots, \underline{v_n} \in V$, definiamo l'insieme delle combinazioni lineari di $\underline{v_1}, \ldots, \underline{v_n}$ come
    $$
        \text{Span}(\underline{v_1}, \ldots, \underline{v_n}) = \left\{ \underline{v} \in V \, \bigg| \, \exists \alpha_1, \ldots, \alpha_n \in \mathbb{K} : \underline{v} = \alpha_1 \underline{v_1} + \ldots + \alpha_n \underline{v_n} \right\}
    $$
\end{definition}
\begin{prop}
    Sia $(V, \mathbb{K}, +, \cdot)$ uno spazio vettoriale sul campo $\mathbb{K}$ e siano $\underline{v_1}, \ldots, \underline{v_n} \in V$. Allora $\text{Span}(\underline{v_1}, \ldots, \underline{v_n})$ è un sottospazio vettoriale di $V$.
\end{prop}
\begin{prop}
    Sia $(V, \mathbb{K}, +, \cdot)$ uno spazio vettoriale sul campo $\mathbb{K}$ e siano $\underline{v_1}, \ldots, \underline{v_n} \in V$. Allora
    $$
        \text{Span}(\underline{v_1}, \ldots, \underline{v_n}) \text{ è il più piccolo sottospazio vettoriale di } V \text{ che contiene } \underline{v_1}, \ldots, \underline{v_n}.
    $$
\end{prop}
\begin{exercise}
    \label{ex:1.2} Mostrare le due proposizioni qua sopra. 
\end{exercise}
Possiamo generalizzare la definizione di $\text{Span}$ di vettori ad insiemi di vettori arbitrari, ovvero
\begin{definition}[Span di un insieme]
    Sia $(V, \mathbb{K}, +, \cdot)$ uno spazio vettoriale sul campo $\mathbb{K}$ e sia $A \subseteq V$ un insieme, allora definiamo lo \text{Span} di $A$ come
    $$
        \text{Span}(A) = \left\{ \underline{v} \in V \, \bigg| \, \exists \underline{v_1}, \ldots, \underline{v_n} \in A : \exists \alpha_1, \ldots, \alpha_n \in \mathbb{K} : \underline{v} = \alpha_1 \underline{v_1} + \ldots \alpha_n \underline{v_n} \right\}.
    $$
\end{definition}
\begin{prop}[Span di un insieme è un sottospazio vettoriale]
    Sia $(V, \mathbb{K}, +,\cdot)$ uno spazio vettoriale sul campo $\mathbb{K}$ e sia $A \subseteq V$ un insieme. Allora $\text{Span}(A)$ è un sottospazio vettoriale.
\end{prop}
\begin{prop}
    Sia $(V, \mathbb{K}, +, \cdot)$ uno spazio vettoriale sul campo $\mathbb{K}$. Se $A \subseteq V$ un sottospazio vettoriale di $V$, allora
    $$
        \text{Span}(A) = A.
    $$
\end{prop}
\begin{prop}
    Sia $(V, \mathbb{K}, +, \cdot)$ uno spazio vettoriale sul campo $\mathbb{K}$. Siano $A, B \subseteq V$ insiemi, allora
    $$
        A \subseteq \text{Span}(B) \implies \text{Span}(A) \subseteq \text{Span}(B).
    $$
\end{prop}
\begin{prop}
    Sia $(V, \mathbb{K}, +, \cdot)$ uno spazio vettoriale sul campo $\mathbb{K}$. Siano $A, B \subseteq V$ insiemi, allora
    $$
        A \subseteq B \implies \text{Span}(A) \subseteq \text{Span}(B).
    $$
\end{prop}
\begin{exercise}
    \label{ex:1.3} Mostrare le quattro proposizioni qua sopra.
\end{exercise}
\section{Vettori linearmente indipedenti}
Un concetto fondamentale che ci aiuterà a capire meglio gli spazi vettoriali è il concetto di \textbf{vettori linearmente indipendenti}, che possiamo definire come segue:
\begin{definition}[vettori linearmente indipendenti]
    Sia $(V, \mathbb{K}, +, \cdot)$ uno spazio vettoriale sul campo $\mathbb{K}$ e siano $\underline{v_1}, \ldots, \underline{v_n} \in V$. Diremo che i vettori $\underline{v_1}, \ldots, \underline{v_n}$ sono linearmente indipendenti se l'equazione
    $$
    \alpha_1 \underline{v_1} + \ldots + \alpha_n \underline{v_n} = \underline{0}
    $$
    ha come \textbf{unica} soluzione quella banale, ovvero $\forall \, i \in \{ 1, \ldots, n \}, \alpha_i = 0$
\end{definition}
Perché siamo interessati a lavorare con i vettori linearmente indipendenti? La motivazione è semplice: le coordinate rispetto a tali vettori sono univoche, come mostriamo adesso
\begin{prop}
    Sia $(V, \mathbb{K}, +, \cdot)$ uno spazio vettoriale sul campo $\mathbb{K}$ e siano $\underline{v_1}, \ldots, \underline{v_n}$ dei vettori linearmente indipedenti, allora se si ha che
    $$
        \alpha_1 \underline{v_1} + \ldots + \alpha_n \underline{v_n} = \alpha_1' \underline{v_1} + \ldots + \alpha_n' \underline{v_n} \implies \forall i \in \{ 1,\ldots, n \}, \, \alpha_i = \alpha_i'
    $$
\end{prop}
\begin{proof}
    Portando tutti i termini a sinistra otteniamo una relazione equivalente,
    $$
        \alpha_1 \underline{v_1} + \ldots + \alpha_n \underline{v_n} = \alpha_1' \underline{v_1} + \ldots + \alpha_n' \underline{v_n} \iff (\alpha_1 - \alpha_1')\underline{v_1} + \ldots + (\alpha_n - \alpha_n') \underline{v_n} = \underline{0},
    $$
    ma, per ipotesi, abbiamo che i vettori $\underline{v_1}, \ldots, \underline{v_n}$ sono linearmente indipendenti, dunque l'unica soluzione a quest'ultima equazione è quella banale, ovvero deve risultare che
    $$
        \forall i \in \{ 1, \ldots, n \}, \alpha_i - \alpha_i' = 0 \implies \forall i \in \{ 1, \ldots, n \}, \alpha_i = \alpha_i',
    $$
    ovvero la tesi.
\end{proof}
\begin{remark}
    Essere vettori linearmente indipedenti \textbf{non} è una proprietà dei singoli vettori, ma una proprietà che hanno i vettori come insieme. Consiglio di guardare il prossimo esempio per chiarire meglio questo concetto
\end{remark}
\begin{example}[$\mathbb{R}^3$ "chiarificatore" di tutto il creato]
    Fissiamo $V = \mathbb{R}^3$ e consideriamo i vettori $\underline{v_1}$
\end{example}